% =====================================================================
%  Quy chuẩn so sánh công bằng các phương pháp pruning
%  *** Biên dịch bằng XeLaTeX / LuaLaTeX *** (Overleaf: Compiler -> XeLaTeX)
% =====================================================================
\documentclass[11pt,a4paper]{article}

\usepackage{fontspec}
\setmainfont{Latin Modern Roman}

\usepackage[margin=2.2cm]{geometry}
\usepackage{array}
\usepackage{tabularx}
\usepackage{amsmath}
\usepackage{enumitem}
\usepackage{hyperref}
\hypersetup{colorlinks=true, linkcolor=blue!50!black, urlcolor=blue!50!black}

\setlist{leftmargin=1.4em,topsep=1pt,itemsep=1pt}
\newcommand{\code}[1]{\texttt{#1}}
\newcolumntype{Y}{>{\raggedright\arraybackslash}X}

\title{Quy chuẩn so sánh công bằng các phương pháp pruning}
\author{}
\date{}

\begin{document}
\maketitle
\vspace{-2.5em}

% =====================================================================
\section{Đặt vấn đề}

So các phương pháp pruning ở cùng target sparsity / compress rate là không công bằng: rate là núm vặn,
không phải thành quả nén. Mỗi phương pháp áp ``rate = 0.5'' khác nhau (uniform / auto / decompose) và
cắt ở vị trí khác nhau. Thành quả nén đo bằng FLOPs, Params, latency, không phải bằng ``sparsity''.

FLOPs và Params không tuyến tính:
\[
  \text{Params} = C_{out}C_{in}k_h k_w, \qquad \text{FLOPs} = \text{Params}\times(H\times W).
\]
$H\times W$ khác nhau mỗi layer nên cùng rate cho ra Params/FLOPs khác nhau, và không khớp đồng thời cả
hai được. Ví dụ: bi-level @rate 0.5 $\to$ 100M FLOPs, mAP 0.33; SVP @rate 0.5 $\to$ 130M FLOPs, mAP
0.34. Ở cùng rate SVP ``thắng'', nhưng chỉ vì nén ít hơn; ở cùng 100M FLOPs bi-level mới thắng.

% =====================================================================
\section{Cách nén của từng phương pháp và yêu cầu trung thành gốc}

Mỗi phương pháp gồm: núm ngân sách, phân bổ per-layer, tiêu chí chọn filter. Ba phần này phải trung
thành bản gốc; chỉ cơ chế cắt/harness được dùng chung.

\footnotesize
\renewcommand{\arraystretch}{1.35}
\begin{tabularx}{\textwidth}{|p{2cm}|Y|Y|Y|p{2.4cm}|}
\hline
Method & Núm = \% của & Phân bổ per-layer & Chọn filter & Trung thành gốc\\
\hline
bi-level & filter/layer & uniform & Song Han $\to$ L1-$\infty$-$\infty$ & (pp của ta)\\
\hline
CORING & filter (tích luỹ) & uniform (gốc: list) & HOSVD $+$ VBD $+$ min\_sum & lịch khớp; đổi allocation về per-layer\\
\hline
NORTON & filter/layer $+$ rank & list per-layer & subspace-angle & rebuild thiếu bước prune\\
\hline
SLIMING & kênh toàn cục & auto (GAM theo SV) & GEM (max nuclear norm) & khớp gốc; thêm đo \%FLOPs\\
\hline
Torch-Pruning & \code{pruning\_ratio} & auto (DepGraph) & magnitude/FPGM/BN & dùng API thư viện\\
\hline
\end{tabularx}
\normalsize

\medskip
Cần sửa trước khi so:
\begin{itemize}
  \item CORING: đổi phân bổ về per-layer list như \code{main\_kshot.py} gốc (hiện uniform).
  \item NORTON: thêm stage 2 prune subspace-angle (port \code{pruning/saliency.py} $+$ \code{prune.py}); hiện chỉ CP-decompose.
  \item SLIMING: giữ GAM (đã khớp); thêm đo \%FLOPs (đang theo \%kênh).
  \item Mọi method: tái hiện số gốc trên benchmark gốc một lần để xác nhận cài đúng.
\end{itemize}

% =====================================================================
\section{Các bài A* so sánh thế nào}

Đọc code + bảng SLIMING/CORING/NORTON/Torch-Pruning: họ không so ở target sparsity mà so trên
\%MACs/\%Params đo được. Bảng SLIMING (Table 3--7) mỗi dòng báo \code{Method | Top-1 | MACs (CR\%) |
Params (CR\%)}; \code{CR} là compression ratio đo được. Cách so:
\begin{itemize}
  \item Nhóm thành tầng compression (vd $\sim$50\%, $\sim$70\% MACs), trong mỗi tầng so accuracy.
  \item Pareto: thắng khi accuracy cao hơn ở compression tương đương hoặc tốt hơn. Ví dụ nguyên văn: ``ARPruning comparable param reduction (89 vs 92\%) but lags accuracy (92.69 vs 93.74)''.
  \item Torch-Pruning cũng đo MACs/Params bằng \code{count\_ops\_and\_params} rồi so accuracy-vs-MACs.
\end{itemize}
Kết luận: trục so là \%MACs/\%Params đo được, theo Pareto, trình bày dạng bảng phân tầng hoặc đường
cong. Target sparsity chỉ là núm nội bộ.

% =====================================================================
\section{Quy chuẩn của chúng ta}

Mốc so:
\begin{itemize}
  \item Trục chính: \%MACs (FLOPs) giảm. Phụ: \%Params, latency/FPS.
  \item NORTON bắt buộc đo latency (CP-decompose depthwise, FLOPs không phản ánh tốc độ).
  \item Chỉ khớp một trục (MACs); trục kia report xấp xỉ (không tuyến tính nên không ép cả hai).
\end{itemize}

Mốc tương đối, tự scale theo model (không hardcode giá trị tuyệt đối):
\[
  R \in \{30,50,70\}\%\ \text{giảm MACs}, \qquad \text{target\_MACs} = (1-R)\times \text{MACs}_0.
\]

% =====================================================================
\section{Run case}

Đo \%Params/\%MACs không cần train (chỉ phụ thuộc cấu trúc đã cắt) --- tách khỏi bước finetune.

\renewcommand{\arraystretch}{1.35}
\begin{tabularx}{\textwidth}{|p{1.7cm}|Y|}
\hline
Phase 0 & Đo dense một lần: $(\text{MACs}_0, \text{Params}_0)$. Cố định cho mọi method: dataset (train/val/test), input, finetune budget, scope.\\
\hline
Phase 1 & Sweep cấu trúc (không train): mỗi (method, rate) $\to$ \code{prune} $\to$ \code{measure} (params, flops). Rẻ, vài giây/điểm.\\
\hline
Phase 2 & Finetune (cùng budget) $\to$ \code{eval} mAP trên test $+$ \code{benchmark} latency $\to$ ghi CSV.\\
\hline
\end{tabularx}

\medskip
Cách rẻ nhất --- sweep $+$ nội suy: chạy grid rate cố định (4--5 mức) $\to$ đường cong mAP-vs-\%MACs
$\to$ đọc mọi target bằng nội suy, không finetune lại. Số finetune $=$ số rate (vd 5 method $\times$ 5
rate $=$ 25 run, vài phút/run). Binary-search chỉ khi cần đúng một mức.

% =====================================================================
\section{Logic code}

\begin{enumerate}
  \item \code{measure(model)} $\to$ (params, flops) qua thop $+$ numel; \code{benchmark\_latency(model)} $\to$ ms/FPS (warm-up $+$ trung bình).
  \item \code{prune(dense, method, rate)} --- interface thống nhất cho 5 method.
  \item \code{compare\_methods.py} --- loop Phase 1--2, ghi CSV \code{(model, method, rate, params, flops, latency, mAP)}.
  \item \code{analyze.py} --- đọc CSV, vẽ đường mAP-vs-\%MACs, nội suy tại grid \{30,50,70\}\%, xuất bảng phân tầng.
\end{enumerate}

Checklist: cùng dense, cùng finetune budget, cùng test set, đo Params/MACs/latency thật, mỗi method
dùng allocation tự nhiên (không ép uniform), so ở iso-\%MACs.

% =====================================================================
\section{So sánh cái gì}

So trên mặt phẳng accuracy vs \%MACs theo Pareto (giống A*):
\begin{itemize}
  \item Bảng phân tầng (kiểu SLIMING Table 3): mỗi tầng compression so mAP các method.
  \item Đường cong mAP-vs-\%MACs: method có đường cao và lệch trái thắng.
  \item Đột phá của bi-level $=$ đường cong trùm CORING/SLIMING/NORTON ở mọi mức \%MACs.
\end{itemize}

Bảng đích (per model, mốc tương đối), kèm cột \%Params $+$ latency; NORTON để bảng riêng có latency:
\begin{center}
\begin{tabular}{@{}lcccc@{}}
\hline
Giảm MACs & bi-level & CORING & SLIMING & NORTON\\
\hline
$-30\%$ & mAP & mAP & mAP & mAP (+lat)\\
$-50\%$ & mAP & mAP & mAP & mAP (+lat)\\
$-70\%$ & mAP & mAP & mAP & mAP (+lat)\\
\hline
\end{tabular}
\end{center}

Tóm: mỗi method nén trung thành gốc $\to$ sweep rate $\to$ đo \%MACs/\%Params/latency/mAP (finetune
cùng budget) $\to$ so ở iso-\%MACs theo Pareto. Rate chỉ để rải điểm; trục so là \%MACs đo được.

\end{document}